\documentclass[11pt]{article}

% Required packages
\usepackage{graphicx}
\usepackage{longtable}
\usepackage{booktabs}
\usepackage{multirow}
\usepackage{hyperref}
\usepackage{array}
\usepackage{caption}
\usepackage{amsmath}

% Set margins (adjust as needed)
\usepackage[margin=1in]{geometry}

\begin{document}

\hfill\break
\textbf{Proceedings of the ASME 2025}

\textbf{International Design Engineering Technical Conferences and\\
Computers and Information in Engineering Conference}

\textbf{IDETC/CIE2025\\
August 17-20, 2025, Anaheim, California}

DETC2025-168841

Developing multi-modal machine learning model for PREDICTING PERFORMANCE
OF AUTOMOTIVE HOOD FRAMEs

{\def\LTcaptype{none} % do not increment counter
\begin{longtable}[]{@{}
  >{\raggedright\arraybackslash}p{(\linewidth - 2\tabcolsep) * \real{0.2366}}
  >{\raggedright\arraybackslash}p{(\linewidth - 2\tabcolsep) * \real{0.2626}}@{}}
\toprule\noalign{}
\begin{minipage}[b]{\linewidth}\centering
\textbf{Abhishek Indupally}

Clemson University

Clemson, SC
\end{minipage} & \begin{minipage}[b]{\linewidth}\centering
\textbf{Satchit Ramnath*}

Clemson University

Clemson, SC
\end{minipage} \\
\midrule\noalign{}
\endhead
\bottomrule\noalign{}
\endlastfoot
\end{longtable}
}

Abstract

\emph{Is there a way for a designer to evaluate the performance of a
given hood frame geometry without spending significant time on
simulation setup? This paper seeks to address this challenge by
developing a multimodal machine-learning (MMML) architecture that learns
from different modalities of the same data, to predict performance
metrics. It also aims to use the MMML architecture to enhance the
efficiency of engineering design processes by reducing reliance on
computationally expensive simulations in the initial phases of design.
The proposed architecture helps accelerate design exploration, enabling
rapid iteration while maintaining high accuracy in predictions. The
study also presents results that show that by combining multiple data
modalities, MMML outperforms traditional single-modality (or unimodal)
approaches. Two new frame geometries, not part of the training dataset,
are also used for prediction using the trained MMML model to showcase
the ability to generalize to unseen frame models. The findings
underscore MMML\textquotesingle s potential in supplementing traditional
simulation-based workflows, particularly in the conceptual design phase,
and highlight its role in bridging the gap between machine learning and
real-world engineering applications. This research paves the way for the
broader adoption of machine learning techniques in engineering design,
with a focus on refining multimodal approaches to optimize structural
development and accelerate the design cycle.}

Keywords: multimodal machine learning, computer-aided design,
computer-aided engineering, geometric dataset, data-driven design,
structural predictions using ML

\begin{enumerate}
\def\labelenumi{\arabic{enumi}.}
\item
  INTRODUCTION
\end{enumerate}

In engineering design, a structured and systematic approach is essential
for developing effective solutions {[}1{]}. One of the primary tasks in
the conceptual design phase is assessing the structural integrity and
overall performance of the design. Traditionally, evaluating structural
integrity requires computational simulations such as Finite Element
Analysis (FEA), which involves defining boundary conditions, running
simulations, and validating results {[}2{]}. While highly accurate, and
also necessary, FEA is computationally expensive and time-intensive,
potentially creating bottlenecks in the iterative design process during
the initial phases. This slows down innovation and increases the design
cycle time, affecting the number of design alternatives explored.
Reducing design cycle time is, therefore, crucial, as it enables
engineers to assess multiple configurations more efficiently, refine
their designs in real time, and accelerate product development without
compromising accuracy.

Machine learning (ML) advancements offer to supplement traditional
simulation-driven design workflows. ML-based models can learn from
existing simulation data and provide fast, approximate performance
metrics predictions, allowing engineers to evaluate design iterations in
real time. This rapid feedback loop enhances the design process by
enabling quick assessments and iterative refinements.

In the context of an automotive hood design, the hood consists of two
primary components: the outer skin and the hood frame, as shown in
\hyperref[_Ref193023952]{Figure 1}. The hood skin refers to the outer,
visible panel (class-A surface) of the hood, defining the
vehicle\textquotesingle s aesthetics and aerodynamics, while the hood
frame is the structural support underneath that panel, providing
structural integrity, crash safety, and attachment points. The frame
consists of various geometrical features, as shown in
\hyperref[_Ref193044630]{Figure 2}. The ribs are considered the primary
features since they provide structural stiffness, and the pockets are
the secondary features that help reduce the weight. Since the hood skin
dictates the design space and the curvature, the frame must be
engineered to follow the same design guidelines while still satisfying
performance requirements such as weight and strength.

\includegraphics[width=1.5in,height=1.14028in]{media/image1.png}
\includegraphics[width=1.5in,height=1.11319in]{media/image2.png}

(a) (b)

\begin{figure}
\centering
\includegraphics[width=3.5in,height=1.70412in,alt={A close-up of a car\textquotesingle s front bumper AI-generated content may be incorrect.}]{media/image3.png}
\caption{\protect\phantomsection\label{_Ref193023952}{}Figure 1: (a)
Automotive hood skin and (b) the frame {[}3{]}}
\end{figure}

\protect\phantomsection\label{_Ref193044630}{}Figure 2: Primary features
on a hood frame (Ribs), secondary features on a hood frame (Pockets)

A designer developing a hood frame must assess its structural integrity
using FEA by predicting key performance metrics, which may include
stresses and deflections. However, relying solely on FEA for these
evaluations is effort and time-intensive, as each iteration requires
re-running simulations, creating a computational bottleneck that limits
rapid design exploration. Effective feedback during the design cycle
{[}4{]} is crucial, as it allows designers to quickly estimate
performance metrics, assess structural integrity, and refine their
concepts without the time-intensive burden of full-scale simulations. To
enable this, a method for quick performance estimation is needed---one
that enables designers to analyze localized effects of geometric
changes, add or remove structural features, and understand how these
modifications impact performance without the computational overhead of
full-scale simulations.

Although these methods aid in conceptual automotive body designs,
predicting structural performance accurately during the design process
requires a model that can effectively capture and interpret complex
geometric relationships. Capturing the complex details and relationships
may not explicitly be available in a dataset, and hence requires new
model architectures to extract, learn, and build the implicit
relationships. This can be achieved by extracting data in different
representations called ``\emph{modalities}''. The concept of modality is
foundational in multimodal machine learning, referring to a specific
representation of information in a particular format. Liang et al.
{[}5{]} describe modality as "a way in which a natural phenomenon is
perceived or expressed". This definition encompasses various types of
representations, including visual, auditory, textual, and numerical
data. Although research has been done to evaluate various ML models to
predict the performance of designs, especially for automotive body
design, the input data to these models are unimodal - only one modality
- like images or point clouds or parameters.

This research proposes a Multi-Modal ML (MMML) architecture that can
process different modalities to make accurate predictions. In this
research, the three primary modalities considered are: image data,
geometric data, and parametric data. The approach integrates image-based
features extracted from convolutional neural networks (CNNs) with
geometric and parametric data to improve the overall predictive
accuracy. \hyperref[_Ref193056771]{Figure 3} shows the workflow for the
MMML model. The images provide spatial and visual information; the
geometric data includes cross-sectional profiles of the design, and the
parametric data includes parameters corresponding to geometric features
in the model. Combining multiple modalities enables a more comprehensive
understanding of a design's behavior. The contributions of this research
are:

\begin{itemize}
\item
  Development of a Custom MMML Architecture -- Design and develop a MMML
  architecture that integrates various modalities to predict key
  performance metrics.
\item
  Comparison with Single-Modality Models -- The proposed MMML
  architecture is evaluated against a baseline model trained with a
  single modality to assess the benefits of multimodal learning.
\item
  Generalization to Unseen Designs -- The model is tested on new hood
  frame designs to evaluate its robustness, predictive accuracy, and
  ability to generalize beyond the training data.
\end{itemize}

\begin{figure}
\centering
\includegraphics[width=3.5in,height=3.11389in]{media/image4.png}
\caption{\protect\phantomsection\label{_Ref193056771}{}Figure 3:
Multi-modal machine learning flowchart}
\end{figure}

\begin{enumerate}
\def\labelenumi{\arabic{enumi}.}
\setcounter{enumi}{1}
\item
  \textbf{BACKGROUND}
\end{enumerate}

Machine learning algorithms are increasingly used in automotive body
design to optimize safety, performance, material efficiency, and cost.
By analyzing large datasets, these algorithms help predict outcomes and
improve decision-making in areas like crashworthiness, panel rigidity,
and material usage. ML can also help automate design tasks, speeding up
development cycles and aiding innovation.

\textbf{2.1 Machine Learning in Automotive Design}

There has been quite a bit of research into integrating ML algorithms
into the domain of automotive design. Liu et al. {[}6{]} proposed a
structured framework to develop design-specific knowledge graphs, which
are then used in deep learning to learn graph embeddings, make
predictions, and support reasoning. Bodendorf and Franke {[}7{]}
compared various machine learning approaches to predict the costs in an
early product design phase of automotive components. The study looked at
the performance of six machine learning algorithms for estimating costs.
Felix et al. {[}8{]} proposed an approach that mirrors the team
structure of engineers, CAD designers, and quality assurance, with a
Vision Language Model (VLM)-based Multi-Agent System with access to
parametric CAD tooling and documentation. Combining agents for
requirements engineering, CAD engineering, and vision-based quality
assurance, a model is generated automatically from sketches and/ or
textual descriptions. Wong et al. {[}9{]} proposed an end-to-end prompt
evolution design optimization (PREDO) framework contextualized in a
vehicle design scenario that leverages a vision-language model to
penalize impractical car designs synthesized by a generative model. The
research used an evolutionary strategy coupled with an optimization
objective function that comprises a physics-based solver and a
vision-language model for practical or functional guidance in the
generated car designs. Rios et al. {[}10{]} propose and realize a fully
automated evolutionary design optimization framework using Shap-E
{[}11{]} in the context of aerodynamic vehicle optimization. Zhen et al.
{[}12{]} proposed a method for automotive hood design based on machine
learning-assisted structural optimization. The objective of the study
was to optimize the hood frame for pedestrian head impact loads.

\textbf{2.2 Multi-Modal Machine Learning}

Multimodal machine learning tasks emerge when inputs or outputs are
represented in various formats or when different types of atomic units
of information are combined within the same system {[}13{]}. These
models have gained significant traction due to their ability to leverage
complementary information from multiple modalities, leading to more
informed predictions. For instance, visual information can be enhanced
with textual descriptions or numerical features to generate more precise
models that account for real-world complexities. Integrating multiple
modalities has proven highly effective in improving prediction accuracy
and generalization across various machine-learning applications.
Baltrusaitis et al. {[}14{]} provide a comprehensive overview of
multimodal machine learning, highlighting five main technical
challenges: representation, translation, alignment, fusion, and
co-learning. Their work emphasizes that multimodal models often
outperform unimodal counterparts by capturing diverse aspects of input
data, thereby mitigating the limitations of any single modality.

Multimodal machine learning (MMML) has demonstrated significant
advancements in predictive modeling, particularly in automotive
engineering. Su et al. {[}15{]} highlight the benefits of integrating
multiple data modalities over unimodal approaches, which rely on a
single data source and may lead to incomplete insights. The MMML model
used in the research combines images, text, and parametric data to
predict five key vehicle rating scores, achieving a 4\% to 12\%
improvement over unimodal models. This underscores the value of
leveraging diverse data types for more accurate predictions. While the
work focuses on vehicle ratings, a similar challenge exists in
evaluating the performance of CAD-based mechanical components.
Structured parametric data often provides the most informative input,
but non-visual data, such as text, can sometimes outperform images in
predictive tasks.

In computer-aided design (CAD), particularly in conceptual design,
integrating multiple data representations can enhance the design
process. Li et al. {[}16{]} introduced LLM4CAD, a large language model
that combines images, sketches, and textual descriptions to generate 3D
geometry. The research underscores the importance of multimodal
approaches in CAD, demonstrating how the fusion of different input
modalities can lead to more expressive and effective design workflows.
By bridging visual and textual representations, MMML facilitates
interactive and sophisticated design processes, enabling designers to
explore and refine concepts more efficiently. A key takeaway from
LLM4CAD is the crucial role of high-quality datasets in the success of
multimodal models. As research in this domain progresses, developing
diverse, domain-specific datasets will be essential for improving
predictive performance and generalizability in engineering design
models.

A study by Gong et al. {[}17{]} on bimodal learning in materials science
highlights the advantages of combining composition and structure data to
predict material properties. The framework processes composition as
vectors and structure as graphs, fusing these modalities to enhance
predictive performance. The results demonstrate that bimodal models
consistently outperform unimodal approaches by capturing complementary
information, reinforcing the value of integrating diverse data types.
The research also highlights the importance of effective data
representation for machine learning applications.

The study by Kumar {[}18{]} on geometrical parameter extraction from
flexible mechanical components and assemblies provides critical insights
into the application of machine learning for mechanical performance
prediction. By leveraging the CarHoods10K {[}19{]} dataset, a
comprehensive collection of 10,000 validated automotive hood frame
geometries, the study establishes a strong foundation for data-driven
design evaluation. One of the key contributions of Kumar's work is the
multi-view approach (using images), which demonstrates how various views
of hood frames can enhance structural understanding. The method uses
images to extract features for performance prediction. However, the
study also reveals a fundamental limitation of image-only models---while
they effectively capture visual and shape-based features, they lack
structured, numerical context that could provide deeper insights into
mechanical behavior.

Hwang et al. {[}20{]} present EMMA, a new approach to improving
autonomous driving systems by using artificial intelligence (AI) to
process and understand multiple types of information at once.
Traditional self-driving systems rely on separate components for tasks
like detecting objects, planning movements, and understanding road
layouts. These components work independently and pass information
between them, which can sometimes lead to errors and inefficiencies.
Instead of using separate modules, EMMA takes in all the necessary
data---like camera images, navigation directions, and vehicle
status---and processes it in a single unified AI model. This allows the
system to make driving decisions more efficiently and accurately. One of
the key innovations is the use of Multimodal Large Language Models
(MLLMs), a type of AI that can combine different types of inputs (such
as images and text) into a single understanding of the driving
environment. While EMMA is a promising step toward smarter self-driving
AI, the authors acknowledge some limitations, such as high computational
requirements and challenges in processing depth and distance information
because it does not use LiDAR (a laser-based distance measurement
technology).

Acosta et al. {[}21{]} explore how integrating multiple data
sources---such as medical images, genomics, biosensors, and clinical
records---can enhance artificial intelligence (AI) applications in
healthcare. The authors highlight the limitations of current AI models,
which typically rely on a single data modality and provide only
moment-in-time assessments, whereas real-world medical decisions require
continuous, multimodal inputs. A key contribution of this work is its
discussion of multimodal AI models, which can analyze diverse biomedical
data types, including wearable sensors, genetic and epigenetic markers,
imaging, clinical text, and social determinants of health. By combining
these sources, multimodal AI enables more accurate diagnosis, prognosis,
and treatment planning, bridging gaps in current medical AI
applications.

The proposed research builds upon these foundations, leveraging
multimodal learning to predict the performance of automotive hood
frames. By combining image data with geometric and parametric
information, the research aims to overcome the limitations of unimodal
and multi-view CNN approaches and provide more accurate, robust
predictions for structural performance.

\begin{enumerate}
\def\labelenumi{\arabic{enumi}.}
\setcounter{enumi}{2}
\item
  \textbf{DATASET EXTRACTION AND PREPROCESSING}
\end{enumerate}

The CarHoods10K dataset offers over 10,000 3D mesh automobile hood frame
geometries developed by an industry-standard CAD generation process
{[}19{]}{[}22{]}. The dataset contains highly realistic,
expert-validated designs that are created by generating over 100
parameterized shape variants (from skins and feature patterns) through a
feature-based modeling approach. This technique enables the creation of
diverse designs by adjusting parameters that influence key features like
ribs and cut-outs, which are essential for the structural integrity of
the hoods. Each model in the dataset is available in STL format, along
with design parameters and performance metrics from finite element
analysis (FEA), such as stress, deformation, and mass. These details are
vital for developing advanced computational and ML models to predict
performance in automotive engineering. \hyperref[_Ref193037678]{Figure
4} displays the 100 shape variants in the geometry dataset. A
feature-based design approach was used to generate the large dataset,
details of which can be found in {[}23{]}. The performance data for the
models in the dataset, which include vonMises stress, direction
deflection, and mass, were obtained from FEA as described in {[}24{]}.
The proposed method uses the dataset to extract data -- images,
geometric data, and parametric data -- for training and validating
advanced architectures that will aid designers in evaluating new designs
rapidly.

\begin{figure}
\centering
\includegraphics[width=3.5in,height=2.41184in]{media/image5.png}
\caption{\protect\phantomsection\label{_Ref193037678}{}Figure 4 :
CarHoods10K dataset shape variations}
\end{figure}

\textbf{3.1 Image Dataset (Top and Side Views)}

The top (or normal direction) view of the hood frame is extracted and
used to give the ML model context about the type, number, and position
of various geometrical features with respect to the design space of the
frame, while the side views give context about the curvature and
``flow'' of the hood frame. By combining these perspectives, the ML
model can better understand and extract the object's shape, contents,
and structure. These features, extracted from the images, are essential
for predicting the performance of the hood frame, such as load-bearing
capacity or stress distribution. The images are processed using two
separate Residual Networks (ResNet50) {[}25{]}. ResNet was chosen
because of the efficiency and speed of training the model. ResNet50
makes predictions by categorizing an image, and in this case, the goal
is to train the model to predict the frame\textquotesingle s performance
based on the extracted features. \hyperref[_Ref193045931]{Figure 5}
shows an example of top and side views of a hood frame.

\includegraphics[width=2in,height=1.64787in]{media/image6.png}

\protect\phantomsection\label{_Ref193045931}{}Figure 5: (a) Top view (b)
side view of a hood frame

\textbf{3.2 Geometric and Parametric Data (Cross-sections \& Feature
Depths)}

Given the change in curvature along the hood frame, extracting the
cross-section at two locations -- 25\% and 75\% along the width -- would
help provide additional information about the frame to the algorithm.
This is the geometric data or modality that is part of the MMML
algorithm. The change in curvature of the frame is not ``visible'' or
available from the image data. \hyperref[_Ref193053659]{Figure 6} shows
the locations at which the cross-sections of the frame are extracted.
The cross-sections are generated by generating planes at the two
locations and then using them to slice the frame model. It is automated

\begin{figure}
\centering
\includegraphics[width=2.88056in,height=1.97986in]{media/image7.png}
\caption{\protect\phantomsection\label{_Ref193053659}{}Figure 6: Frame
with the 25\% and 75\% location for cross-sections}
\end{figure}

A custom Python script is used to generate the cross-sections from the
geometry files in the dataset. The extracted cross-sections are recorded
as \((x,z)\) points in CSV files -- the representation type is textual
data. However, due to variations in the size and shape of hood frames in
the dataset, the number of points in a cross-section was not the same.
Hence a standardized representation was established by defining a
minimum required number of coordinates per cross-section to address
this. The first and last 10\% of the points were used to pad the
boundary coordinates to preserve the overall shape and maintain a
consistent length of the input data. This approach ensured consistency
across the dataset while retaining the key geometric features of each
cross-section. \hyperref[_Ref193054282]{Figure 7} shows an example of
the cross-section extracted from a hood frame.

\begin{figure}
\centering
\includegraphics[width=3.21701in,height=2.5in,alt={A graph with a blue line Description automatically generated}]{media/image8.png}
\caption{\protect\phantomsection\label{_Ref193054282}{}Figure 7: Example
of cross-section points extracted from a hood frame}
\end{figure}

Although the cross-sections provide information on the variation in
curvature of the hood frame, the ability to accurately capture
\emph{all} feature depths depends on the location of the slicing plane.
In some cases, certain critical structural features may not intersect
with the slicing plane which may impact the performance predictions. To
address this limitation, the parametric values for depths of the primary
features (or ribs) are extracted from the dataset using an automated
script. This ensures that key geometric characteristics influencing
structural performance are incorporated into the model, even if it is
not captured in the cross-sections. This information is stored in CSV
files.

A total of 20 design points from the final dataset of images,
cross-sections, and parameters were saved separately for validating the
model after training. These points were not used in the training or
testing of the model. The rest of the data was split 80/20 for train and
test datasets. The discussion in the results section utilizes the
validation data to draw conclusions on the performance and efficiency of
the model.

\begin{enumerate}
\def\labelenumi{\arabic{enumi}.}
\setcounter{enumi}{3}
\item
  \textbf{MODEL ARCHITECTURE}
\end{enumerate}

The research proposes to develop a multimodal neural network
architecture to predict the performance metrics of car hood frames by
integrating the three key data modalities: images, geometric data, and
parametric data, described earlier. The proposed architecture consists
of \emph{five} individual networks and one final set of fully connected
layers. \hyperref[_Ref193056345]{Figure 8} shows the MMML architecture
developed in this research. The \emph{five} individual networks are:

\begin{enumerate}
\def\labelenumi{\arabic{enumi}.}
\item
  Two separate convolutional networks (Residual Networks -- ResNet50) to

  \begin{enumerate}
  \def\labelenumii{\alph{enumii}.}
  \item
    Extract feature maps from the top views
  \item
    Extract feature maps from the side views
  \end{enumerate}
\item
  Two neural networks to learn from the two cross-section datasets
\item
  One neural network to process the depth parameters in the frame model
\end{enumerate}

Other networks like VGGNet {[}26{]}, ResNet100, and AlexNet {[}27{]}
were also tested and compared against ResNet50. The accuracy and
computation resources, based on the dataset, remained unchanged. Hence
the ResNet50 model was chosen as the network to process images in the
MMML architecture. The input to each network is separate, and the
networks don't interact with each other. Once all the relevant features
are extracted from all the modalities, the results are concatenated and
pooled into a unified representation. This is then fed into a network of
fully connected layers to make the final prediction on the performance
of the hood frame. By leveraging multimodal learning, this approach
ensures a comprehensive understanding of design-performance
relationships, enhancing prediction accuracy while reducing reliance on
computationally expensive simulations.

\begin{figure}
\centering
\includegraphics[width=3.5in,height=5.48519in]{media/image9.png}
\caption{\protect\phantomsection\label{_Ref193056345}{}Figure 8:
Multi-Modal Machine Learning Architecture}
\end{figure}

\begin{enumerate}
\def\labelenumi{\arabic{enumi}.}
\item
  \textbf{Creating Image Embeddings}
\end{enumerate}

The images are represented as *.jpg files of size 128 X 128 X 1 pixels -
they are scaled and converted to grayscale. Given that a key
characteristic of CNN models is extracting features, it is done by
identifying various components of an image using filters. These filters,
or vectors with assigned weights, are applied to an image to generate
feature maps. The network attempts to learn and extract important
patterns from the input data through this process. From a mathematical
perspective, an image is represented as a matrix, where each cell
corresponds to a pixel, and its value indicates the
pixel\textquotesingle s brightness on a scale from 0 to 1. Each pixel is
assigned a weight, which can be referred to as its intensity. A filter,
or a smaller matrix, is then applied to the image resulting in a feature
map, which is a two-dimensional array that captures essential
information about the features within the image.
\hyperref[_Ref193035853]{Figure 9} shows a sample of feature maps and a
filter for a hood frame model in the dataset. Each feature map
represents an attempt to identify a significant feature in the image.
Given the dataset involves images of automotive hood frames, the CNN is
expected to capture key geometric features, represented visually, such
as the ribs, pockets, surface curvature, etc. For example, a top-view
provides insight into how the hood features (primary and secondary ribs,
pockets, hinges, and locks) are laid out and where the features on the
hood are located. Similarly, a side-view shows the significance of the
depth and curvature of the hood frame. These features will help the
model better understand the shape and size of the features, along with
their relative positions.

\includegraphics[width=1in,height=0.76478in,alt={A close-up of a heat map AI-generated content may be incorrect.}]{media/image10.jpeg}\includegraphics[width=1in,height=0.77014in,alt={A close-up of a screen AI-generated content may be incorrect.}]{media/image11.jpeg}\includegraphics[width=1in,height=0.76669in,alt={A close-up of a face AI-generated content may be incorrect.}]{media/image12.jpeg}

(a)

\includegraphics[width=1in,height=1.26687in,alt={A chart of a graph AI-generated content may be incorrect.}]{media/image13.png}

(b)

\protect\phantomsection\label{_Ref193035853}{}Figure 9: Example
visualization of (a) feature maps \& (b) filter

A ResNet50 base model is initialized without the ``top layers'', also
known as the classification layers. The model is initialized without
pre-trained weights; in other words, the weights are randomly assigned.
This allows the model to learn directly from the dataset rather than
relying on past datasets it has trained on. The extracted features are
processed through a global pooling layer to reduce the spatial
dimensions and produce a feature vector for each image. A fully
connected layer (128 modes deep and a ReLU activation function) followed
by batch normalization is used to capture complex nonlinear
relationships within the data and to normalize activations to enhance
training stability and convergence. As a result, the ResNet50 model
outputs a 128-dimensional feature vector, which effectively encodes the
most significant aspects of the image. \hyperref[_Ref193060321]{Figure
10} shows the architecture of the two image-processing networks.

\begin{figure}
\centering
\includegraphics[width=2.74889in,height=1.35404in,alt={A diagram of a computer AI-generated content may be incorrect.}]{media/image14.png}
\caption{\protect\phantomsection\label{_Ref193060321}{}Figure 10:
ResNet50 architectures for top and side view images}
\end{figure}

The feature vector from the CNN model would have captured all the
important information from the images since, in a CNN model, with each
subsequent convolution layer, the feature maps influence changes, and
they focus on the finer details of the features as opposed to the whole
image. These fine details of the features in the image would then be
able to provide vital information to predict the performance of a hood
frame.

\begin{enumerate}
\def\labelenumi{\arabic{enumi}.}
\setcounter{enumi}{1}
\item
  \textbf{Creating Parametric Data Embedding}
\end{enumerate}

Neural network models were created for processing the data associated
with the two cross-sections and parametric data. The structure of the
networks processing the cross-section data is shown in
\hyperref[_Ref193061552]{Figure 11}. The input layer takes in the
cross-section data, followed by a series of hidden layers. Each network
has a total of six hidden layers of sizes 128, 256, 512, 1024, 512, and
128 nodes each, all utilizing the ReLU activation function.

\includegraphics[width=1.50182in,height=2.3in,alt={A diagram of a number of units AI-generated content may be incorrect.}]{media/image15.png}\includegraphics[width=1.46225in,height=2.3in,alt={A diagram of a number of units AI-generated content may be incorrect.}]{media/image16.png}

\protect\phantomsection\label{_Ref193061552}{}Figure 11: Neural network
architectures for cross-section data

A Batch Normalization layer is applied to stabilize activations and
improve training efficiency, functioning similarly to the batch
normalization layer in the ResNet50 model. Finally, the output layer
generates a 128-dimensional feature vector that embeds and captures the
most critical features of the cross-sections.

The structure of the network processing the parametric data is shown in
\hyperref[_Ref193061796]{Figure 12}. The input layer takes in the
parameter (or depth) data, followed by a series of hidden layers. The
network has a total of four hidden layers of sizes 128, 256, 512, and
128 nodes each, all utilizing the ReLU activation function. Similar to
the networks processing the cross-sections, a batch normalization layer
is used, which then generates a 128-dimensional vector that embeds the
parametric data.

\begin{figure}
\centering
\includegraphics[width=1.57222in,height=1.90566in]{media/image17.png}
\caption{\protect\phantomsection\label{_Ref193061796}{}Figure 12: Neural
network for parametric data}
\end{figure}

\begin{enumerate}
\def\labelenumi{\arabic{enumi}.}
\setcounter{enumi}{2}
\item
  \textbf{Performance Prediction}
\end{enumerate}

The fully connected layers play a crucial role in the final prediction.
Once the significant features are extracted, the fully connected layers
use these embeddings to make predictions on the output parameters. The
fully connected layers thus serve as the decision-making component,
interpreting various aspects of the input data and translating them into
meaningful predictions regarding its performance.

In the context of hood frame prediction, the final component of the
multimodal model is a neural network designed to predict the performance
metrics of car hood frames. The feature embeddings from the ResNet50
models and the three neural networks are concatenated to create a
unified representation. \hyperref[_Ref193062912]{Figure 13} shows the
architecture of the fully connected layers used to make the final
prediction. The concatenated vector is processed through a series of
hidden layers consisting of 128, 64, 32, 16, and 3 nodes each. The final
dense layer includes three nodes, each responsible for predicting a
distinct performance metric. L2 regularization is applied to prevent
overfitting, ensuring that excessively large weight values are penalized
during training. Additionally, dropout layers are utilized, with neurons
being randomly deactivated at rates of 50\%, 40\%, and 30\% across
different layers. This regularization technique reduces reliance on
specific neurons, thereby enhancing generalization. Finally, the output
from the network consists of three units, each predicting a performance
metric for the car hood frames.

\begin{figure}
\centering
\includegraphics[width=1.14653in,height=2.38681in]{media/image18.png}
\caption{\protect\phantomsection\label{_Ref193062912}{}Figure 13: Fully
connected layers for final prediction}
\end{figure}

\begin{enumerate}
\def\labelenumi{\arabic{enumi}.}
\setcounter{enumi}{4}
\item
  \textbf{RESULTS AND DISCUSSION}
\end{enumerate}

The most accurate predictions from a MMML model are achieved by
fine-tuning various model parameters. Optimizing a model's architecture
involves several key adjustments, such as modifying the number of layers
(in both the initial and final models) and updating the number of nodes
in the neural network. One of the main challenges in this process is
striking the right balance between model complexity and simplicity.
Finding the right model architecture and parameter required careful
tuning. A model that was too simple would risk overfitting the training
data, leading to poor generalization on the testing dataset. Conversely,
an overly complex model could potentially struggle to extract meaningful
insights from the data. The architecture described earlier in this paper
represents the best balance tested so far---providing strong predictive
performance without excessive complexity.

In addition to optimizing the model's architecture, fine-tuning the
hyperparameters is equally critical in improving the performance of the
model. Hyperparameters such as learning rate, batch size, patience
value, and the number of training epochs all significantly impact the
model's performance. The learning rate was one of the most important
hyperparameters for the MMML model. The initial learning rate was set to
1e-5, which aids in rapid adjustments, but during training, as the
number of epochs increases, both training and validation loss tend to
plateau. This was countered using a dynamic function that updated (or
lowered) the learning rate every time the loss plateaued. This dynamic
nature of adjusting the learning rate enabled finer adjustments for
improved performance and convergence.

\begin{figure}
\centering
\includegraphics[width=3.40208in,height=2.61321in,alt={A graph of a graph AI-generated content may be incorrect.}]{media/image19.png}
\caption{\protect\phantomsection\label{_Ref193066000}{}Figure 14:
Training and validation loss vs number of epochs}
\end{figure}

The convergence behavior of the training and validation loss over 100
epochs is shown in \hyperref[_Ref193066000]{Figure 14}. The term
validation loss in the figure is actually the ``test'' loss and not from
the validation dataset. The validation dataset is used only after the
training is complete. Initially, both losses exhibit a sharp decline,
indicating rapid learning of key features from the data. A sudden spike
in the validation loss at the beginning is primarily due to using random
weights instead of pre-trained weights. Around the 20-epoch mark, the
validation loss stabilizes, suggesting that the model begins refining
its learned representations rather than simply memorizing patterns.
Beyond 50 epochs, both losses approach an asymptotic limit, with minimal
further reduction, demonstrating that the model has reached an optimal
state. Notably, the validation loss remains closely aligned with the
training loss throughout the training process. This indicates that the
model is neither overfitting nor underfitting but generalizes well to
unseen data. The absence of significant divergence between the two
losses suggests that appropriate regularization techniques were applied,
ensuring stable learning dynamics. The smooth convergence further
highlights the benefits of an adaptive learning rate strategy,
preventing premature stagnation or excessive oscillations in model
performance. The model was trained on a Windows operating system, using
an NVIDIA RTX A5000 GPU running on a 13th Gen Intel Core i9-13900K
processor with 64GB RAM.

\textbf{5.1 Results from Validation Dataset}

Once training was complete, the model was validated using the validation
dataset. Three separate graphs were generated to assess the model
performance. Each plot shows the prediction of a performance metric --
von Mises stress, mass, or deflection. The figures below illustrate the
model's average error across these categories and individual error
distributions for each data point.

\includegraphics[width=3.56in,height=1.99945in,alt={A graph of a graph with numbers and a line AI-generated content may be incorrect.}]{media/image20.png}

\begin{enumerate}
\def\labelenumi{(\alph{enumi})}
\item
  vM stress (MPa) vs design points
\end{enumerate}

\includegraphics[width=3.56in,height=2.06998in,alt={A graph with numbers and lines AI-generated content may be incorrect.}]{media/image21.png}

\begin{enumerate}
\def\labelenumi{(\alph{enumi})}
\setcounter{enumi}{1}
\item
  Mass (kg) vs design points
\end{enumerate}

\includegraphics[width=3.56in,height=2.03627in,alt={A graph with numbers and lines AI-generated content may be incorrect.}]{media/image22.png}

\begin{enumerate}
\def\labelenumi{(\alph{enumi})}
\setcounter{enumi}{2}
\item
  Directional deflection (mm) vs design points
\end{enumerate}

\protect\phantomsection\label{_Ref193068838}{}Figure 15: Actual vs
predicted values for the validation data (a) vM stress, (b) mass, and
(c) directional deflection

\begin{longtable}[]{@{}
  >{\centering\arraybackslash}p{(\linewidth - 2\tabcolsep) * \real{0.1760}}
  >{\centering\arraybackslash}p{(\linewidth - 2\tabcolsep) * \real{0.0988}}@{}}
\caption{Table 1: Average error in the predicted values for the
CarHoods10k dataset (based on \hyperref[_Ref193068838]{Figure
15})}\tabularnewline
\toprule\noalign{}
\begin{minipage}[b]{\linewidth}\centering
\textbf{Output Parameter}
\end{minipage} & \begin{minipage}[b]{\linewidth}\centering
\textbf{Error \%}
\end{minipage} \\
\midrule\noalign{}
\endfirsthead
\toprule\noalign{}
\begin{minipage}[b]{\linewidth}\centering
\textbf{Output Parameter}
\end{minipage} & \begin{minipage}[b]{\linewidth}\centering
\textbf{Error \%}
\end{minipage} \\
\midrule\noalign{}
\endhead
\bottomrule\noalign{}
\endlastfoot
vM Stress & 5.48\% \\
Mass & 9.27\% \\
Deflection & 16.84\% \\
\end{longtable}

The MMML model performed quite well on the validation dataset, with
error percentages well below 20\%. These results were compared to
another model using the same dataset and performance metrics as
discussed by Kumar et al. in {[}18{]}. The results from this research
are significantly better than the unimodal model (trained only on
images) used in the other research. Kumar et al. had a confidence
interval of 20\% with a model accuracy of 90\%, as compared to the MMML
model, which had an accuracy of over 94\% and confidence intervals in
single digits for stress and mass and around 17\% for deflection as
noted in the table above.

5.2 Model Performance on Unseen Hood Frame Designs

To evaluate the generalizability of the proposed model, two hood frame
designs, not part of CarHoods10k, were tested. The designs were sourced
from an online community platform (GrabCAD {[}28{]}).
\hyperref[_Ref193069767]{Figure 16} and \hyperref[_Ref193069768]{Figure
17} show the top and side views, along with the details on the
cross-sections, of the two hood frames. These models underwent the same
preprocessing steps, including data extraction and finite element
analysis (FEA) using a similar model setup (to the CarHoods10k dataset).
The FEA model setup for the two frames is shown in
\hyperref[_Ref193070448]{Figure 18}. The FEA results served as ground
truth, allowing the proposed model to compare the predicted performance
values. \hyperref[_Ref193070777]{Figure 19} shows the model's
predictions of performance compared to the actual FEA results, and
\hyperref[_Ref193070835]{Table 2} shows the average error in prediction.

\begin{longtable}[]{@{}
  >{\centering\arraybackslash}p{(\linewidth - 2\tabcolsep) * \real{0.2373}}
  >{\centering\arraybackslash}p{(\linewidth - 2\tabcolsep) * \real{0.2373}}@{}}
\caption{\protect\phantomsection\label{_Ref193069767}{}Figure 16: (a)
Top view, (b) cross-section locations, (c) side view, and (d)
cross-section points for Model 1}\tabularnewline
\toprule\noalign{}
\begin{minipage}[b]{\linewidth}\centering
\includegraphics[width=1.5in,height=0.78417in]{media/image23.png}
\end{minipage} & \begin{minipage}[b]{\linewidth}\centering
\includegraphics[width=1.5in,height=0.86489in,alt={A drawing of a curved road AI-generated content may be incorrect.}]{media/image24.png}
\end{minipage} \\
\midrule\noalign{}
\endfirsthead
\toprule\noalign{}
\begin{minipage}[b]{\linewidth}\centering
\includegraphics[width=1.5in,height=0.78417in]{media/image23.png}
\end{minipage} & \begin{minipage}[b]{\linewidth}\centering
\includegraphics[width=1.5in,height=0.86489in,alt={A drawing of a curved road AI-generated content may be incorrect.}]{media/image24.png}
\end{minipage} \\
\midrule\noalign{}
\endhead
\bottomrule\noalign{}
\endlastfoot
(a) & (b) \\
\includegraphics[width=1.5in,height=0.35588in,alt={A paper model of a spaceship AI-generated content may be incorrect.}]{media/image25.png}
&
\includegraphics[width=1.49955in,height=0.93913in]{media/image26.png} \\
(c) & (d) \\
\end{longtable}

\begin{longtable}[]{@{}
  >{\centering\arraybackslash}p{(\linewidth - 2\tabcolsep) * \real{0.2373}}
  >{\centering\arraybackslash}p{(\linewidth - 2\tabcolsep) * \real{0.2373}}@{}}
\caption{\protect\phantomsection\label{_Ref193069768}{}Figure 17: (a)
Top view, (b) cross-section locations, (c) side view, and (d)
cross-section points for Model 2}\tabularnewline
\toprule\noalign{}
\begin{minipage}[b]{\linewidth}\centering
\includegraphics[width=1.5in,height=1.51941in,alt={A black object with holes AI-generated content may be incorrect.}]{media/image27.jpeg}
\end{minipage} & \begin{minipage}[b]{\linewidth}\centering
\includegraphics[width=1.5in,height=0.77577in,alt={A drawing of a curved object AI-generated content may be incorrect.}]{media/image28.png}
\end{minipage} \\
\midrule\noalign{}
\endfirsthead
\toprule\noalign{}
\begin{minipage}[b]{\linewidth}\centering
\includegraphics[width=1.5in,height=1.51941in,alt={A black object with holes AI-generated content may be incorrect.}]{media/image27.jpeg}
\end{minipage} & \begin{minipage}[b]{\linewidth}\centering
\includegraphics[width=1.5in,height=0.77577in,alt={A drawing of a curved object AI-generated content may be incorrect.}]{media/image28.png}
\end{minipage} \\
\midrule\noalign{}
\endhead
\bottomrule\noalign{}
\endlastfoot
(a) & (b) \\
\includegraphics[width=1.5in,height=0.25643in,alt={A close up of a roof AI-generated content may be incorrect.}]{media/image29.jpeg}
&
\includegraphics[width=1.49955in,height=0.96522in]{media/image30.png} \\
(c) & (d) \\
\end{longtable}

\begin{longtable}[]{@{}
  >{\centering\arraybackslash}p{(\linewidth - 10\tabcolsep) * \real{0.0917}}
  >{\centering\arraybackslash}p{(\linewidth - 10\tabcolsep) * \real{0.0677}}
  >{\centering\arraybackslash}p{(\linewidth - 10\tabcolsep) * \real{0.0887}}
  >{\centering\arraybackslash}p{(\linewidth - 10\tabcolsep) * \real{0.0677}}
  >{\centering\arraybackslash}p{(\linewidth - 10\tabcolsep) * \real{0.0887}}
  >{\centering\arraybackslash}p{(\linewidth - 10\tabcolsep) * \real{0.0718}}@{}}
\caption{\protect\phantomsection\label{_Ref193070448}{}Figure 18: FEA
setup for hood frames}\tabularnewline
\toprule\noalign{}
\begin{minipage}[b]{\linewidth}\centering
\includegraphics[width=1.5in,height=0.78416in,alt={A close-up of a metal object AI-generated content may be incorrect.}]{media/image31.png}
\end{minipage} & \begin{minipage}[b]{\linewidth}\centering
\includegraphics[width=2in,height=1.39016in,alt={A model of a race track AI-generated content may be incorrect.}]{media/image32.png}
\end{minipage} \\
\midrule\noalign{}
\endfirsthead
\toprule\noalign{}
\begin{minipage}[b]{\linewidth}\centering
\includegraphics[width=1.5in,height=0.78416in,alt={A close-up of a metal object AI-generated content may be incorrect.}]{media/image31.png}
\end{minipage} & \begin{minipage}[b]{\linewidth}\centering
\includegraphics[width=2in,height=1.39016in,alt={A model of a race track AI-generated content may be incorrect.}]{media/image32.png}
\end{minipage} \\
\midrule\noalign{}
\endhead
\bottomrule\noalign{}
\endlastfoot
(a) & (b) \\
\includegraphics[width=1.4in,height=1.42708in]{media/image33.png} &
\includegraphics[width=2.18681in,height=1.42708in]{media/image34.png} \\
(c) & (d) \\
\end{longtable}

As shown in \hyperref[_Ref193070777]{Figure 19}(a), the predicted stress
values are slightly lower than the actual values (from FEA), with an
average underprediction of 18.52\%. While the model tends to be
conservative in its stress estimates, the predictions are still within a
reasonable range, capturing the general trend of stress distribution.
This indicates that the model is not randomly over- or underestimating
values but rather making consistent and physically plausible
predictions, avoiding extreme overfitting to the training data.

\hyperref[_Ref193070777]{Figure 19}(b) shows that the model overpredicts
mass, with an average error of 32.94\%. While there is a tendency to
overestimate mass, this does not indicate overfitting. Instead, the
model slightly biases mass estimates towards safer, heavier designs.
Given that mass estimation is influenced by multiple geometric
parameters, small shifts in feature representation could explain these
variations. Importantly, the model still captures the correct order of
magnitude and does not diverge drastically from the actual value.

\begin{figure}
\centering
\includegraphics[width=3.56042in,height=2.14028in]{media/image35.png}
\caption{(a)}
\end{figure}

\begin{figure}
\centering
\includegraphics[width=3.56042in,height=2.14028in]{media/image37.png}
\caption{\includegraphics[width=3.56042in,height=2.14028in]{media/image36.png}\\
(b)}
\end{figure}

(c)

\protect\phantomsection\label{_Ref193070777}{}Figure 19: Model
prediction vs actual values (a) stress, (b) mass, and (c) deflection

\hyperref[_Ref193070777]{Figure 19}(c) shows that the model's deflection
predictions are lower than the FEA results with an average error of
28.33\%. One potential reason for underprediction could be the presence
of geometrical features (and their shapes) that are unique to this
frame, and nothing remotely similar exists in the training dataset.
However, the predictions remain physically reasonable and trend in the
correct direction, reinforcing that the model is capturing the overall
behavior of hood frame designs rather than memorizing specific patterns.

The results indicate a strong performance of the model, given that the
new hood frames are not a part of the training dataset and the MMML
model has never seen these before. This demonstrates the MMML model's
robustness and ability to generalize beyond the training dataset.

\begin{longtable}[]{@{}
  >{\centering\arraybackslash}p{(\linewidth - 10\tabcolsep) * \real{0.0917}}
  >{\centering\arraybackslash}p{(\linewidth - 10\tabcolsep) * \real{0.0677}}
  >{\centering\arraybackslash}p{(\linewidth - 10\tabcolsep) * \real{0.0887}}
  >{\centering\arraybackslash}p{(\linewidth - 10\tabcolsep) * \real{0.0677}}
  >{\centering\arraybackslash}p{(\linewidth - 10\tabcolsep) * \real{0.0887}}
  >{\centering\arraybackslash}p{(\linewidth - 10\tabcolsep) * \real{0.0718}}@{}}
\caption{\protect\phantomsection\label{_Ref193070835}{}Table 2: Actual
vs predicted values (and average error) for new hood
models}\tabularnewline
\toprule\noalign{}
\begin{minipage}[b]{\linewidth}\raggedright
\textbf{Model}

\textbf{Output}
\end{minipage} &
\multicolumn{2}{>{\centering\arraybackslash}p{(\linewidth - 10\tabcolsep) * \real{0.1564} + 2\tabcolsep}}{%
\begin{minipage}[b]{\linewidth}\centering
\textbf{Hood Frame 1}
\end{minipage}} &
\multicolumn{2}{>{\centering\arraybackslash}p{(\linewidth - 10\tabcolsep) * \real{0.1564} + 2\tabcolsep}}{%
\begin{minipage}[b]{\linewidth}\centering
\textbf{Hood Frame 2}
\end{minipage}} & \begin{minipage}[b]{\linewidth}\centering
\textbf{Avg. Error}
\end{minipage} \\
\midrule\noalign{}
\endfirsthead
\toprule\noalign{}
\begin{minipage}[b]{\linewidth}\raggedright
\textbf{Model}

\textbf{Output}
\end{minipage} &
\multicolumn{2}{>{\centering\arraybackslash}p{(\linewidth - 10\tabcolsep) * \real{0.1564} + 2\tabcolsep}}{%
\begin{minipage}[b]{\linewidth}\centering
\textbf{Hood Frame 1}
\end{minipage}} &
\multicolumn{2}{>{\centering\arraybackslash}p{(\linewidth - 10\tabcolsep) * \real{0.1564} + 2\tabcolsep}}{%
\begin{minipage}[b]{\linewidth}\centering
\textbf{Hood Frame 2}
\end{minipage}} & \begin{minipage}[b]{\linewidth}\centering
\textbf{Avg. Error}
\end{minipage} \\
\midrule\noalign{}
\endhead
\bottomrule\noalign{}
\endlastfoot
& \textbf{Actual Value} & \textbf{Predicted Value} & \textbf{Actual
Value} & \textbf{Predicted Value} & \textbf{\%} \\
\textbf{vM Stress (MPa)} & 207.5 & 156.88 & 152.49 & 133.19 & 18.52\% \\
\textbf{Geometry Mass (kg)} & 9.68 & 13.24 & 10.53 & 13.60 & 32.94\% \\
\textbf{Dir. Def. (mm)} & 10.03 & 7.78 & 11.74 & 7.78 & 28.33\% \\
\end{longtable}

5.3 Comparing Unimodal vs Multimodal Performance

To quantitatively assess the effectiveness of the multimodal approach,
the predictions from the MMML model were compared to predictions from a
unimodal (image-based) architecture. Although the unimodal architecture
only used images, it used a multiview-CNN approach {[}29{]}, to consider
both the top and side views in making the predictions. The evaluation
was conducted for both the CarHoods10K datasets, as well as the two new
hood frames.

\includegraphics[width=2.75in,height=2.20242in,alt={A graph showing different values AI-generated content may be incorrect.}]{media/image38.jpeg}

(a)

\includegraphics[width=2.75in,height=2.24179in,alt={A graph showing different values AI-generated content may be incorrect.}]{media/image39.jpeg}

(b)

\includegraphics[width=2.75in,height=2.25051in,alt={A graph of different values AI-generated content may be incorrect.}]{media/image40.jpeg}

(c)

\protect\phantomsection\label{_Ref193073344}{}Figure 20: Average model
performance for unimodal (images only) vs multimodal for CarHoods10k and
new hood frames

The (average) prediction accuracies for the two models are compared to
the actual values (from FEA) shown in \hyperref[_Ref193073344]{Figure
20}. The stress prediction results (\hyperref[_Ref193073344]{Figure
20}(a)) demonstrate that the multimodal model consistently outperforms
the unimodal model by providing predictions closer to the actual values.
The results are summarized in \hyperref[_Ref193074567]{Table 3}.

\begin{longtable}[]{@{}
  >{\centering\arraybackslash}p{(\linewidth - 8\tabcolsep) * \real{0.1235}}
  >{\centering\arraybackslash}p{(\linewidth - 8\tabcolsep) * \real{0.1079}}
  >{\centering\arraybackslash}p{(\linewidth - 8\tabcolsep) * \real{0.0816}}
  >{\centering\arraybackslash}p{(\linewidth - 8\tabcolsep) * \real{0.0816}}
  >{\centering\arraybackslash}p{(\linewidth - 8\tabcolsep) * \real{0.0812}}@{}}
\caption{\protect\phantomsection\label{_Ref193074567}{}Table 3: Summary
of average error \% in prediction for the two models on two different
datasets}\tabularnewline
\toprule\noalign{}
\begin{minipage}[b]{\linewidth}\centering
Dataset
\end{minipage} & \begin{minipage}[b]{\linewidth}\centering
Model
\end{minipage} & \begin{minipage}[b]{\linewidth}\centering
Stress Error \%
\end{minipage} & \begin{minipage}[b]{\linewidth}\centering
Mass Error \%
\end{minipage} & \begin{minipage}[b]{\linewidth}\centering
Def. Error \%
\end{minipage} \\
\midrule\noalign{}
\endfirsthead
\toprule\noalign{}
\begin{minipage}[b]{\linewidth}\centering
Dataset
\end{minipage} & \begin{minipage}[b]{\linewidth}\centering
Model
\end{minipage} & \begin{minipage}[b]{\linewidth}\centering
Stress Error \%
\end{minipage} & \begin{minipage}[b]{\linewidth}\centering
Mass Error \%
\end{minipage} & \begin{minipage}[b]{\linewidth}\centering
Def. Error \%
\end{minipage} \\
\midrule\noalign{}
\endhead
\bottomrule\noalign{}
\endlastfoot
\multirow{2}{=}{\centering\arraybackslash CarHoods10k} & Unimodal &
8.2\% & 8.1\% & 15.2\% \\
& Multimodal & 3.6\% & 4.7\% & 10.5\% \\
\multirow{2}{=}{\centering\arraybackslash New Models} & Unimodal &
50.3\% & 9.6\% & 48\% \\
& Multimodal & 19.4\% & 32.9\% & 28.7\% \\
\end{longtable}

The results indicate that the MMML model consistently performs better on
all the performance metrics on the CarHoods10k dataset. It also performs
extremely well in improving the predictions of the stress and deflection
values for the two unseen hood frame models. Overall, these findings
demonstrate that multimodal models leveraging geometric, parametric, and
image-based information outperform unimodal approaches, especially in
out-of-distribution scenarios. The ability to maintain predictive
accuracy on a previously unseen dataset supports the argument that
multimodal learning is a robust framework for structural performance
prediction tasks.

\textbf{5.4 Summary}

These results indicate that the multimodal learning model successfully
generalizes to new designs without signs of overfitting (which would
manifest as near-perfect predictions for training data but poor
performance on unseen data) or severe underfitting (resulting in
completely random or highly inaccurate estimates). The stress and
deflection predictions are slightly conservative, which is preferable in
many engineering applications where safety is a priority. While slightly
overestimated, the mass predictions remain within a reasonable deviation
and do not indicate an erratic or unstructured response. Importantly,
the model correctly captures the relative differences between the two
new designs, reinforcing that it has captured meaningful structural
relationships.

\begin{enumerate}
\def\labelenumi{\arabic{enumi}.}
\setcounter{enumi}{5}
\item
  \textbf{CONCLUSION}
\end{enumerate}

By reducing reliance on computationally expensive simulations, the
proposed MMML architecture enhances the efficiency of design
exploration, enabling engineers to iterate rapidly while maintaining
high-performance standards. This work contributes to the broader
adoption of machine learning in engineering design, demonstrating its
potential to supplement traditional simulation-based workflows and
accelerate the development of optimized structures, especially in the
conceptual design phase.

The proposed architecture also demonstrates that MMML outperforms
single-modality approaches, highlighting the benefit of combining and
using \emph{different forms of the same data}. Additionally, the MMML
model was tested on unseen frame models, showcasing its potential for
generalization beyond the training dataset. The key contributions in the
work can be summarized as:

\begin{itemize}
\item
  Developed an MMML architecture and demonstrated the use of multiple
  modalities to effectively train ML models
\item
  Compared the performance of a multimodal architecture to an unimodal
  architecture and validated the effect of utilizing multiple modalities
  to yield better results
\item
  Showcased the model's ability to extend beyond the CarHoods10K
  dataset, demonstrating its generalization to unseen data. However,
  additional testing on varied hood frame designs is necessary for
  broader applicability.
\end{itemize}

Ultimately, this study highlights the potential of MMML in reducing
reliance on simulation, accelerating the design cycle, and enabling
rapid performance evaluation. As machine learning continues to integrate
into engineering workflows, refining multimodal approaches will be key
to bridging the gap between predictive models and real-world
applications.

\textbf{6.1 Future Work}

Despite these strengths, several areas for improvement remain. The
current parametric representation, which relies on cross-sectional data,
may not fully capture intricate geometric features such as ribs or
pockets, which play a crucial role in structural performance. Future
work will explore alternative parametric representations, potentially
incorporating more granular feature-based data and geometric parameters
to provide additional insight into the geometries. Similarly, the image
input could be refined to focus on extracted features rather than full
geometry images, allowing the model to learn more relevant patterns
without redundant information. Future iterations of this model should
investigate more structured feature selection techniques to preserve
essential design characteristics while maintaining computational
efficiency.

As the framework continues to evolve, balancing prediction accuracy with
generalizability remains a priority. While reducing error rates is
desirable, overly precise models risk overfitting, limiting their
applicability to new and diverse designs. By continuously testing on
significantly different hood frame geometries, future research will
assess the robustness of the model across a broader range of engineering
applications.

\textbf{ACKNOWLEDGEMENTS}

The authors would like to thank Simon Vaglienti and Sandesh Kumawat for
their contributions to the project.\\
\strut \\
\textbf{REFERENCES}

{[}1{]} Pahl, G., and Beitz, W., 1996, \emph{Engineering Design---A
System- Atic Approach}, Springer London.

{[}2{]} Heisig, C. P., and Williams-Harkcom, M. A., 1987, ``Integration
of CAD and FEA for Design of Fabricated Structures.''
https://doi.org/10.4271/872017.

{[}3{]} Kim, D.-H., Jung, K.-H., Kim, D.-J., Park, S.-H., Kim, D.-H.,
Lim, J., Nam, B.-G., and Kim, H.-S., 2017, ``Improving Pedestrian Safety
via the Optimization of Composite Hood Structures for Automobiles Based
on the Equivalent Static Load Method,'' Compos Struct, \textbf{176}, pp.
780--789. https://doi.org/10.1016/j.compstruct.2017.06.016.

{[}4{]} Wynn, D. C., and Maier, A. M., 2022, ``Feedback Systems in the
Design and Development Process,'' Res Eng Des, \textbf{33}(3), pp.
273--306. https://doi.org/10.1007/s00163-022-00386-z.

{[}5{]} Liang, P. P., Zadeh, A., and Morency, L.-P., 2022, ``Foundations
and Trends in Multimodal Machine Learning: Principles, Challenges, and
Open Questions,'' ArXiv.

{[}6{]} Liu, A., Zhang, D., Wang, Y., and Xu, X., 2022, ``Knowledge
Graph with Machine Learning for Product Design,'' CIRP Annals,
\textbf{71}(1), pp. 117--120.
https://doi.org/10.1016/j.cirp.2022.03.025.

{[}7{]} Bodendorf, F., and Franke, J., 2021, ``A Machine Learning
Approach to Estimate Product Costs in the Early Product Design Phase: A
Use Case from the Automotive Industry,'' Procedia CIRP, \textbf{100},
pp. 643--648. https://doi.org/10.1016/j.procir.2021.05.137.

{[}8{]} Ocker, F., Menzel, S., Sadik, A., and Rios, T., 2025, ``From
Idea to CAD: A Language Model-Driven Multi-Agent System for
Collaborative Design.''

{[}9{]} Wong, M., Rios, T., Menzel, S., and Ong, Y. S., 2024, ``Prompt
Evolutionary Design Optimization with Generative Shape and
Vision-Language Models,'' \emph{2024 IEEE Congress on Evolutionary
Computation (CEC)}, IEEE, pp. 1--8.
https://doi.org/10.1109/CEC60901.2024.10611898.

{[}10{]} Rios, T., Menzel, S., and Sendhoff, B., 2023, ``Large Language
and Text-to-3D Models for Engineering Design Optimization,'' \emph{2023
IEEE Symposium Series on Computational Intelligence (SSCI)}, IEEE, pp.
1704--1711. https://doi.org/10.1109/SSCI52147.2023.10371898.

{[}11{]} Jun, H., and Nichol, A., 2023, ``Shap-E: Generating Conditional
3D Implicit Functions.''

{[}12{]} Zhan, Z., Fengyao, L., Xin, R., Zhou, G., Zhao, S., He, X.,
Wang, J., and Li, J., 2023, ``Automotive Hood Design Based on Machine
Learning and Structural Design Optimization.''
https://doi.org/10.4271/2023-01-0744.

{[}13{]} Parcalabescu, L., Trost, N., and Frank, A., 2021, ``What Is
Multimodality?''

{[}14{]} Baltrušaitis, T., Ahuja, C., and Morency, L.-P., 2017,
``Multimodal Machine Learning: A Survey and Taxonomy.''

{[}15{]} Su, H., Song, B., and Ahmed, F., 2023, ``Multi-Modal Machine
Learning for Vehicle Rating Predictions Using Image, Text, and
Parametric Data,'' \emph{Volume 2: 43rd Computers and Information in
Engineering Conference (CIE)}, American Society of Mechanical Engineers.
https://doi.org/10.1115/DETC2023-115076.

{[}16{]} Li, X., Sun, Y., and Sha, Z., 2024, ``LLM4CAD: Multi-Modal
Large Language Models for 3D Computer-Aided Design Generation,''
\emph{Volume 6: 36th International Conference on Design Theory and
Methodology (DTM)}, American Society of Mechanical Engineers.
https://doi.org/10.1115/DETC2024-143740.

{[}17{]} Gong, S., Wang, S., Zhu, T., Shao-Horn, Y., and Grossman, J.
C., 2023, ``Multimodal Machine Learning for Materials Science:
Composition-Structure Bimodal Learning for Experimentally Measured
Properties.''

{[}18{]} Kumar, P., 2024, ``A Study on the Extraction of Geometrical
Parameters from Flexible Mechanical Components and Assemblies and Their
Impact on Performance: A Machine Learning Approach,'' M.S. Thesis,
Arizona State University. {[}Online{]}. Available:
https://keep.lib.asu.edu/items/195239. {[}Accessed: 15-Mar-2025{]}.

{[}19{]} Wollstadt, P., Bujny, M., Ramnath, S., Shah, J. J., Detwiler,
D., and Menzel, S., 2022, ``CarHoods10k: An Industry-Grade Data Set for
Representation Learning and Design Optimization in Engineering
Applications,'' IEEE Transactions on Evolutionary Computation, pp. 1--1.
https://doi.org/10.1109/TEVC.2022.3147013.

{[}20{]} Hwang, J.-J., Xu, R., Lin, H., Hung, W.-C., Ji, J., Choi, K.,
Huang, D., He, T., Covington, P., Sapp, B., Zhou, Y., Guo, J., Anguelov,
D., and Tan, M., 2024, ``EMMA: End-to-End Multimodal Model for
Autonomous Driving.''

{[}21{]} Acosta, J. N., Falcone, G. J., Rajpurkar, P., and Topol, E. J.,
2022, ``Multimodal Biomedical AI,'' Nat Med, \textbf{28}(9), pp.
1773--1784. https://doi.org/10.1038/s41591-022-01981-2.

{[}22{]} Ramnath, S., Shah, J. J., Wollstadt, P., Bujny, M., Menzel, S.,
and Detwiler, D., 2022, ``OSU-Honda Automobile Hood Dataset
(CarHoods10k),'' \emph{Dryad Dataset}.
https://doi.org/10.5061/dryad.2fqz612pt.

{[}23{]} Ramnath, S., Haghighi, P., Kim, J. H., Detwiler, D., Berry, M.,
Shah, J. J., Aulig, N., Wollstadt, P., and Menzel, S., 2019,
``Automatically Generating 60,000 CAD Variants for Big Data
Applications,'' \emph{International Design Engineering Technical
Conferences and Computers and Information in Engineering Conference}.

{[}24{]} Ramnath, S., Haghighi, P., Ma, J., Shah, J. J., and Detwiler,
D., 2020, ``Design Science Meets Data Science: Curating Large Design
Datasets for Engineered Artifacts,'' \emph{International Design
Engineering Technical Conferences and Computers and Information in
Engineering Conference}.

{[}25{]} He, K., Zhang, X., Ren, S., and Sun, J., 2016, ``Deep Residual
Learning for Image Recognition,'' \emph{Proceedings of the IEEE Computer
Society Conference on Computer Vision and Pattern Recognition}, pp.
770--778. https://doi.org/10.1109/CVPR.2016.90.

{[}26{]} Simonyan Karen, and Zisserman Andrew, 2015, ``Very Deep
Convolutional Networks for Large-Scale Image Recognition,'' \emph{3rd
International Conference on Learning Representations, ICLR 2015 -
Conference Track Proceedings}, p. 14. {[}Online{]}. Available:
https://www.researchgate.net/publication/265385906\_Very\_Deep\_Convolutional\_Networks\_for\_Large-Scale\_Image\_Recognition\%0Ahttp://www.robots.ox.ac.uk/.

{[}27{]} Krizhevsky, A., Sutskever, I., and Hinton, G. E., 2017,
``ImageNet Classification with Deep Convolutional Neural Networks,''
Commun ACM, \textbf{60}(6), pp. 84--90. https://doi.org/10.1145/3065386.

{[}28{]} GrabCAD, ``GrabCAD Library.'' {[}Online{]}. Available:
https://grabcad.com/. {[}Accessed: 15-Mar-2025{]}.

{[}29{]} Su, H., Maji, S., Kalogerakis, E., and Learned-Miller, E.,
2015, ``Multi-View Convolutional Neural Networks for 3D Shape
Recognition.''

~
